%% friday_experiments.tex  –– v2 (corrected)
%% Preliminary ("Friday") Experiment Proposals
%% Free-surface flows of shear-thickening suspensions
%% APS / RevTeX 4-2 preprint format
%%
%% C. Burton, Driscoll Lab, Northwestern University
%% May 2026
%%
%% v2 corrections:
%%   Exp 1 – fixed H_c derivation (3D Newtonian uses axisymmetric Huppert
%%            scaling, H_c is V-dependent; Saffman-Taylor analogy clarified)
%%   Exp 2 – removed oscillatory "forced oobleck wave" protocol (amplitude
%%            analysis shows v_c requires unreachable stroke at oobleck
%%            frequencies); corrected tau_c numerical estimate; focus on
%%            step-forcing protocol only
%%   Exp 3 – replaced fringe-projection shadowgraphy (Cobelli 2009) with
%%            side-view LED backlit imaging; removed cross-sectional PIV
%%            of opaque cornstarch as primary deliverable (physically
%%            impossible without index-matched system); velocity-reversal
%%            measurement demoted to future extension
%%   Exp 4 – removed claim lambda -> 2pi/k_max as Omega -> 0 (not derived
%%            from Balmforth 2025); Omega* retained as order-of-magnitude only
%%   Exp 5 – corrected b* scaling: delta ~ (2 l tau_c / rho g)^{1/2}, not
%%            the 1/3-power formula; corrected numerical estimate

\documentclass[aps,preprint,superscriptaddress,longbibliography]{revtex4-2}

\usepackage{amsmath}
\usepackage{amssymb}
\usepackage{graphicx}
\usepackage{xcolor}
\usepackage{hyperref}
\usepackage{booktabs}
\usepackage{enumitem}
\usepackage[framemethod=TikZ]{mdframed}

%% ---------------------------------------------------------------
\mdfdefinestyle{questionbox}{
  backgroundcolor=blue!7, linecolor=blue!50, linewidth=1pt,
  roundcorner=4pt, innertopmargin=6pt, innerbottommargin=6pt,
  innerleftmargin=8pt, innerrightmargin=8pt,
  skipabove=6pt, skipbelow=4pt}
\mdfdefinestyle{setupbox}{
  backgroundcolor=green!7, linecolor=green!50, linewidth=1pt,
  roundcorner=4pt, innertopmargin=6pt, innerbottommargin=6pt,
  innerleftmargin=8pt, innerrightmargin=8pt,
  skipabove=6pt, skipbelow=4pt}
\mdfdefinestyle{noveltybox}{
  backgroundcolor=orange!8, linecolor=orange!60, linewidth=1pt,
  roundcorner=4pt, innertopmargin=6pt, innerbottommargin=6pt,
  innerleftmargin=8pt, innerrightmargin=8pt,
  skipabove=6pt, skipbelow=4pt}

\setlength{\emergencystretch}{3em}
\hyphenation{quasi-two-di-men-sion-al}

\newcommand{\ULD}{{U^{\mathrm{LD}}}}
\newcommand{\tauc}{\tau_c}
\newcommand{\etaz}{\eta_0}
\newcommand{\phione}{\phi_1}
\newcommand{\phizero}{\phi_0}

%% ---------------------------------------------------------------
\begin{document}

\title{Friday Experiments: Novel Free-Surface Configurations for Probing
Shear-Thickening Suspension Dynamics}

\author{C.~Burton}
\affiliation{Department of Mechanical Engineering, Northwestern University,
Evanston, IL 60208, USA}

\date{\today}

\begin{abstract}
We propose five preliminary experimental programs---``Friday experiments''---
motivated by two recent landmark papers on shear-thickening (``oobleck'') flows:
Darbois~Texier \textit{et al.}~\cite{DarboisTexier2020,DarboisTexier2023},
who identified and characterized the inertia-free surface-wave instability
of dense cornstarch suspensions on inclined planes, and Bougouin \textit{et al.}\
\cite{Bougouin2026PRL}, who discovered the ``liquid dam''---a constant-speed
gravity-driven front arising from a jammed, load-bearing structure at the
advancing edge.  Each proposed experiment targets a gap that is clearly open in
the literature and feasibly addressable with existing high-speed photography and
a modest laser-based velocimetry system.  The five programs are:
(1)~a three-dimensional radial free-surface dam break that removes the
quasi-two-dimensional channel constraint of all prior liquid-dam work and
probes azimuthal stability of the jammed front; (2)~upward elastic-membrane
forcing of a resting suspension layer (``reverse impact''), providing a
controlled, inertia-free route to jamming-front nucleation and propagation;
(3)~quantitative mapping of the oobleck-wave dispersion relation via controlled
sinusoidal inlet perturbations, providing the first systematic experimental test
of the linear stability theory using only side-view imaging;
(4)~stairwave wavelength selection under a programmable tilt-rate protocol; and
(5)~liquid-dam interaction with a geometric step discontinuity, testing the
intrinsic character of $\ULD$ and providing the first direct measurement of the
overflow-layer thickness $\delta$.  For each experiment we derive the relevant
scalings, identify the decisive observables, and discuss measurement strategy.
\end{abstract}

\maketitle

%% ---------------------------------------------------------------
\section{Introduction and Literature Landscape}
\label{sec:intro}

Dense suspensions of hard particles near the jamming transition exhibit
a remarkably rich class of free-surface flow phenomena absent in Newtonian or
conventional non-Newtonian fluids.  The physical origin lies in the Wyart--Cates
(WC) model of discontinuous shear thickening (DST): a smooth crossover from a
frictionless viscosity $\etaz(\phi) = \eta_s (\phizero - \phi)^{-2}$ at low
stress to a frictional, near-divergent viscosity at stresses above a
concentration-dependent critical stress $\tauc(\phi)$, with a fully re-entrant
(backward-bending) constitutive curve in the shear-jamming (SJ) regime
$\phione < \phi < \phizero$~\cite{WyartCates2014}.

\paragraph{Oobleck waves.}
Darbois~Texier \textit{et al.}~\cite{DarboisTexier2020} reported that dense
cornstarch suspensions flowing down an inclined channel spontaneously develop
prominent, regular surface waves at volume fractions $\phi \gtrsim 0.38$.
Unlike the classical Kapitza (roll-wave) instability of Newtonian thin films,
which requires a finite Reynolds number, oobleck waves appear at vanishingly
small Re and are entirely inertia-free.  Their origin is the coupling between
the free surface and the backward-bending branch of the constitutive curve:
an increase in surface stress reduces flow velocity, which raises the local
free-surface height, which further amplifies the stress---a self-sustaining
feedback~\cite{DarboisTexier2023,Depoilly2024}.  Darbois~Texier
\textit{et al.}~\cite{DarboisTexier2023} subsequently provided a quantitative
linear stability theory spanning the full Kapitza-to-Oobleck range, and
Balmforth~\cite{Balmforth2025} extended this to nonlinear saturation, predicting
that finite-amplitude waves steepen into staircase fronts---precisely the
``stairwaves'' subsequently observed by Bougouin \textit{et al.}~\cite{Bougouin2026PRL}.

\paragraph{Liquid dam.}
Bougouin \textit{et al.}~\cite{Bougouin2026PRL} showed that when a DST
suspension in the SJ regime is released in a dam-break geometry, it spreads
not by decelerating like any known fluid but instead with a nearly vertical
front---a ``liquid dam''---that advances at a constant speed
\begin{equation}
  \ULD(\phi) = 2k\,\frac{\tauc(\phi)^2}{\rho g\,\etaz(\phi)},
  \label{eq:ULD}
\end{equation}
independent of released volume, flow height, and channel slope.  The mechanism
is a jammed, frictional front of length $\ell \sim k H / k_0$ (where $k = 10$
and $k_0 \approx 2$--$3$ from their model, consistent with $\ell/H \approx 3$--$5$
observed in their surface velocimetry) above which a thin frictionless overflow
layer of thickness $\delta$ carries the mass flux.  Equation~\eqref{eq:ULD}
is validated across nearly two decades in speed by varying $\phi$, solvent
viscosity, and particle type.  Related localized jammed structures---frictional
solitons---had previously been identified in pipe flow~\cite{Bougouin2024PNAS}
and Couette cells~\cite{Rathee2022PNAS}, establishing a generic picture of
stress-localization in DST systems.

\paragraph{Open questions.}
The liquid-dam experiments are all performed in quasi-two-dimensional channels;
the three-dimensional geometry is completely unexplored.  The overflow-layer
thickness $\delta$ and jammed-front length $\ell$ appear in the model but have
not been directly measured; the paper explicitly calls for future internal imaging
\cite{Bougouin2026PRL}.  The linear stability theory for oobleck waves has been
tested only against naturally arising waves; no controlled-perturbation experiment
has mapped the dispersion relation $\omega(k)$ or growth rate $\sigma(k)$.  The
stairwave wavelength formula $\lambda \approx 2h_\text{ini}/\sin\theta$ has not
been tested against a controlled tilt protocol.  And the behavior of the liquid
dam upon encountering a geometric discontinuity is entirely unexplored.  The five
experiments proposed below address these gaps one by one.

Throughout, we characterize the system by: particle volume fraction $\phi$
(corrected for cornstarch porosity via $\phi_\text{eff} = 1.136\,\phi_\text{nom}$
\cite{HanVanHaJaeger2017}), critical stress $\tauc(\phi)$ and frictionless
viscosity $\etaz(\phi)$ from WC fits to Couette rheometry (representative values
for cornstarch-water in the SJ regime: $\tauc \approx 2$--$4$~Pa,
$\etaz \approx 0.2$--$1.5$~Pa$\cdot$s~\cite{Bougouin2026PRL}), layer thickness
$H$, jammed-front length $\ell \approx 4H$ (from surface velocimetry of
Bougouin \textit{et al.}), and suspension density
$\rho \approx \rho_w\left[1 + (\rho_p/\rho_w - 1)\phi\right]$.

%% ---------------------------------------------------------------
\section{Experiment~1: Three-Dimensional Radial Free-Surface Dam Break}
\label{sec:3Ddambreak}

\subsection{Physical motivation and gap}

Every existing liquid-dam experiment confines the suspension between parallel
channel walls, reducing the problem to one spatial dimension.  Real
gravity-driven flows of DST suspensions---lava analogues, slurry spills,
concrete placement---are three-dimensional.  When the advancing front is a
growing annulus (radius $R(t)$), azimuthal perturbations of the jammed ring
experience a qualitatively different geometry than the flat channel case.  In
standard viscous spreading, the 3D (axisymmetric) Newtonian gravity current
% CORRECTED 2026-07-06: fixed-VOLUME axisymmetric current is R ~ t^{1/8} (Huppert 1982);
% t^{1/2} is the constant-FLUX case. Planar fixed-volume height decays as t^{-1/5},
% not t^{-2/5}. Applied here and at the three sites below (+ summary list).
follows $R \propto t^{1/8}$ (Huppert 1982~\cite{Huppert1982}), with a
qualitatively different height evolution than the 2D case ($X \propto t^{1/5}$).
Whether the liquid-dam regime persists in 3D, and how the transition to
Newtonian spreading is modified, are both unexplored.

A further question is azimuthal stability.  In the channel, the jammed front
is a rectangle driven forward by the frictionless pool.  In the radial geometry,
the frictionless pool (viscosity $\etaz$) drives the jammed annulus (effective
viscosity $\gg \etaz$) outward.  This configuration---a less-viscous fluid
displacing a more-viscous one in radial outward spreading---is the same
orientation as the Saffman--Taylor instability in a Hele-Shaw cell, raising
the possibility that fingers of the frictionless pool penetrate the jammed ring.
Whether this analogy carries over quantitatively is unknown: the jammed front
is not a passive fluid but a self-sustaining frictional structure.  Determining
whether the front fingers or remains azimuthally uniform is a key open question.

\begin{mdframed}[style=questionbox]
\textbf{Key questions.}
(1)~Does $\ULD = dR/dt$ remain constant and consistent with
Eq.~\eqref{eq:ULD} in radial 3D spreading?
(2)~Is the liquid dam azimuthally uniform, or do fingers of the frictionless
pool penetrate the jammed annulus?
(3)~How does the height evolution $H(t)$ differ from the 2D case, and where
does the transition to Newtonian axisymmetric spreading
($R \propto t^{1/8}$) occur?
\end{mdframed}

\subsection{Geometry and scaling}

Consider a cylindrical release of volume $V = \pi R_0^2 H_0$ spreading
radially on a horizontal surface.  If the flat-top liquid-dam morphology holds
in 3D (as it does in 2D), volume conservation with $h(r,t) \approx H(t)$ for
$r < R(t)$ gives
\begin{equation}
  H(t) = \frac{V}{\pi R(t)^2}.
  \label{eq:3D_conservation}
\end{equation}
The front mass balance is local: a cross-section of the annular front at any
azimuthal position still satisfies $H\,\ULD \sim \delta\bar{u}$ (the same
relation as in 2D), provided the overflow remains frictionless and $\ell \sim H$.
If these conditions hold, the front speed in 3D is still $\ULD(\phi)$ from
Eq.~\eqref{eq:ULD}, giving
\begin{equation}
  R(t) \approx \ULD(\phi)\,t, \qquad
  H(t) \approx \frac{V}{\pi\,\ULD^2\,t^2}.
  \label{eq:3D_kinematics}
\end{equation}
The $H \sim t^{-2}$ decay in 3D is qualitatively faster than the $H \sim
t^{-1/5}$ decay in the 2D Newtonian case and represents a clean, falsifiable
difference from any existing spreading model.

The transition to Newtonian spreading occurs when the Newtonian axisymmetric
front speed overtakes $\ULD$.  For the axisymmetric viscous gravity current
(Huppert 1982 \cite{Huppert1982}), the lubrication scaling gives
$\dot{R}_\text{Newt} \sim \rho g H^3/(\eta_0 R)$, so the transition condition
$\ULD = \dot{R}_\text{Newt}$ yields
\begin{equation}
  \ULD \sim \frac{\rho g H_c^3}{\etaz R_c},
  \label{eq:3D_transition}
\end{equation}
where $R_c$ and $H_c$ are simultaneously constrained by
$H_c = V/(\pi R_c^2)$.  Unlike the 2D case---where $H_c$ depends only on
$\phi$ and $A$---in 3D the transition height and radius depend on $V$ as well.
The coupled system of Eq.~\eqref{eq:3D_transition} and volume conservation
must be solved numerically for each $V$ and $\phi$; subsequent spreading after
the transition follows $R \propto t^{1/8}$ (Huppert 1982 \cite{Huppert1982}).

\begin{mdframed}[style=setupbox]
\textbf{Proposed setup.}
A Plexiglas cylinder ($R_0 = 4$--$6$~cm, $H_0 = 3$--$8$~cm) is filled with
a DST cornstarch suspension at $\phi \in \{0.43, 0.45, 0.46\}$ and placed on
a flat, hydrophobic-coated acrylic sheet.  The cylinder is lifted rapidly
($\Delta t \lesssim 0.1\,H_0/\ULD \sim 0.1$~s) by a pneumatic actuator,
releasing the suspension to spread freely.  Top-view high-speed photography
(1000--3000~fps) records $R(\theta, t)$ where $\theta$ is the azimuthal angle.
A synchronized side-view camera at 500--1000~fps tracks $H(t)$ and verifies
the flat-top morphology.  Surface PTV: the free surface is seeded with
$50$--$100$~\textmu m hollow glass microspheres; a horizontal CW 532~nm
laser sheet illuminates the surface, and the top-view camera (with long-pass
filter) records particle tracks at 1000--2000~fps for radial velocity mapping.
\end{mdframed}

\subsection{Observables and expected results}

From top-view imagery: mean radius $\bar{R}(t)$ tests Eq.~\eqref{eq:3D_kinematics};
a log-log plot of $\bar{R}$ vs.\ $t$ should show slope 1 during the liquid-dam
phase, transitioning to slope $1/2$ after $t_c$.  The rms azimuthal fluctuation
$\sigma_R(t) = \langle [R(\theta,t) - \bar{R}]^2\rangle^{1/2}$ quantifies any
fingering instability.  If fingers grow, their count $m$ and growth rate can be
extracted by azimuthal Fourier decomposition.

From surface PTV: the radial velocity profile $u_r(r,t)$ should exhibit a
localized overspeed near the front (analogous to the 2D signature of Bougouin
\textit{et al.}~\cite{Bougouin2026PRL}) with azimuthal modulation if fingers
are present.

\begin{mdframed}[style=noveltybox]
\textbf{Novelty and publishability.}
This is the first liquid-dam experiment without lateral confinement.  It tests
whether Eq.~\eqref{eq:ULD} is a genuinely local, geometry-independent
constitutive law, reveals whether the jammed annulus is azimuthally stable,
and provides the first experimental measurement of the 3D height evolution
$H(t) \sim t^{-2}$.  A clean result---stability or instability---is a natural
\textit{Physical Review Letters} paper; a full dataset mapping $R_c$, $H_c$,
and $R(\theta, t)$ as functions of $V$ and $\phi$ suits
\textit{Physical Review Fluids}.
\end{mdframed}

%% ---------------------------------------------------------------
\section{Experiment~2: Elastic Membrane Upward Forcing (``Reverse Impact'')}
\label{sec:elasticmem}

\subsection{Physical motivation and gap}

Work from the Driscoll laboratory on drop impact of DST suspensions onto
compliant (PEGDA/PEGMEA gel) substrates has established that the elastic wave
reflection timescale $t_r = 2h/c_s$ competes with jamming front propagation
and strongly alters impact morphology~\cite{Mobaseri2025PNAS,Mobaseri2026arXiv}.
In all of this work, the suspension is the impactor and the elastic substrate
is the target; the jammed region nucleates at the impact interface and propagates
upward through the suspension.

Here we propose to invert this geometry: a thin elastic membrane (stretched
latex or PDMS sheet) is clamped over a rigid frame and loaded from above with
a layer of DST suspension of height $h_0$.  A rod pushes upward through the
frame at controlled velocity $v_\text{rod}$.  There is no impact---the forcing
is slow and controlled---so inertia plays no role.  The free surface of the
suspension (the top surface) is unconfined.  This provides a clean way to study
jamming-front nucleation and propagation in the quasi-static limit, which is
inaccessible in drop-impact experiments where inertia and stress localization
are entangled.

\begin{mdframed}[style=questionbox]
\textbf{Key questions.}
(1)~At what critical velocity $v_c$ does a jammed patch nucleate above the rod
contact, and does $v_c \sim \tauc h_0 / \etaz$ (the WC-based threshold)?
(2)~Does a radially propagating jamming front appear above this threshold,
and how does its speed depend on $v_\text{rod}$ and $\phi$?
(3)~How does the elastic restoring force of the membrane modify the spatial
extent of the jammed region compared to the rigid-substrate impact case?
(4)~Does the free surface deform upward (``dome'') above the jammed patch,
and at what height does the dome arrest?
\end{mdframed}

\subsection{Geometry and scaling}

When the rod tip (radius $a_\text{rod}$) pushes up at speed $v_\text{rod}$,
the suspension layer is compressed vertically at a strain rate
$\dot{\varepsilon} \approx v_\text{rod}/h_0$.  For a suspension layer with
free-slip at the top surface and no-slip at the membrane, the resulting shear
rate is of the same order, $\dot{\gamma} \sim v_\text{rod}/h_0$.  Jamming
nucleates when the local stress exceeds $\tauc$:
\begin{equation}
  v_\text{rod} \geq v_c \equiv \frac{\tauc(\phi)\,h_0}{\etaz(\phi)}.
  \label{eq:vc}
\end{equation}
For cornstarch-water in the SJ regime ($\tauc \sim 2$--$4$~Pa,
$\etaz \sim 0.2$--$1.5$~Pa$\cdot$s) with $h_0 = 1$~cm, this gives
$v_c \sim 10$--$100$~mm/s, accessible with a standard lead-screw actuator or
pneumatic piston.  The elastic membrane contributes a spatially varying upward
normal stress $\sigma_\text{mem}(r) \approx T\,\nabla^2 w(r)$ (where $T$ is
membrane tension and $w$ is the deflection profile), which partially offsets
the rod forcing away from the contact point.  As a result, the jammed region
may arrest at a finite radius $R_\text{jam}$ before the rod is withdrawn,
allowing one to map $R_\text{jam}(v_\text{rod}, T, \phi)$.

Once a jammed front nucleates, its radial propagation speed $c_f$ in the
drop-impact geometry has been shown to scale with the impact velocity and
current front radius~\cite{WaitukaitisJaeger2012,Rathee2022PNAS}.  Whether
the same scaling holds in the quasi-static, upward-forcing geometry---where
neither the forcing velocity nor the stress distribution are the same as in
impact---is an open question.

\begin{mdframed}[style=setupbox]
\textbf{Proposed setup.}
A circular latex membrane ($R_\text{frame} = 10$~cm, thickness 0.3--1~mm,
pretension $T$ set by a calibrated tightening ring; $T$ measured via the
fundamental drum frequency with a microphone) is clamped over a rigid acrylic
support with a central hole (diameter $2a_\text{frame} = 1$~cm).  A suspension
layer ($h_0 = 0.5$--$2$~cm, $\phi = 0.43$--$0.46$) is loaded on top.  A
rounded Delrin rod ($a_\text{rod} = 2$--$5$~mm tip radius) is driven from below
by a stepper-motor lead screw at controlled $v_\text{rod} = 1$--$200$~mm/s.
Top-view high-speed camera (500--2000~fps) images plan-view jamming front
growth.  Side-view camera records free-surface dome profile $h(r,t)$.  Surface
PTV (seeded free surface, horizontal CW laser sheet, top-view camera with
long-pass filter) provides radial velocity $u_r(r,t)$.
\end{mdframed}

\subsection{Observables and expected results}

Vary $v_\text{rod}$ from well below to well above $v_c(\phi)$ at fixed $h_0$
and $\phi$: the onset of visible jamming (central dark rigid patch in top-view)
defines $v_c$ experimentally.  Plotting $v_c$ vs.\ $\tauc h_0/\etaz$ across
multiple $\phi$ tests Eq.~\eqref{eq:vc}.  Above threshold: track $R_\text{jam}(t)$
to extract the front propagation speed $c_f(v_\text{rod}, \phi)$ and compare
with the drop-impact scaling~\cite{WaitukaitisJaeger2012}.  Vary $T$ at fixed
$v_\text{rod}$ and $\phi$: if membrane tension halts front growth at
$R_\text{jam} < R_\text{frame}$, this directly reveals the length scale over
which the elastic restoring stress becomes comparable to $\tauc$.  The dome
height $h_\text{dome}(R_\text{jam}, T, \phi)$ from the side-view camera provides
a complementary observable.

\begin{mdframed}[style=noveltybox]
\textbf{Novelty and publishability.}
This is the first experiment in which an elastic substrate drives jamming
in a resting suspension from below.  By removing impact inertia entirely, it
isolates the quasi-static stress-localization mechanism and provides a direct,
parameter-free test of the WC nucleation threshold $v_c$.  The dependence of
$R_\text{jam}$ on membrane tension $T$ introduces a new length scale with no
analogue in the rigid-substrate impact literature, and connects directly to the
Driscoll-lab program on compliance effects in DST impact.  Target:
\textit{Physical Review Letters} or \textit{Soft Matter}.
\end{mdframed}

%% ---------------------------------------------------------------
\section{Experiment~3: Oobleck-Wave Dispersion Relation via Controlled
Inlet Perturbation}
\label{sec:dispersion}

\subsection{Physical motivation and gap}

Darbois~Texier \textit{et al.}~\cite{DarboisTexier2023} derived the full linear
stability problem for a WC suspension on an inclined plane, predicting the
spatial growth rate $\sigma(k;\,\phi,\theta)$ and phase speed $c(k;\,\phi,\theta)$
as functions of streamwise wavenumber $k$.  Depoilly \textit{et al.}\
\cite{Depoilly2024} generalized this to an arbitrary generalized-Newtonian
rheology, unifying the Kapitza and Oobleck branches.  These theories make
quantitative, parameter-free predictions (all inputs come from rheometry): the
neutral stability curve $k_\text{neutral}(\phi,\theta)$, the wavenumber of
maximum growth $k_\text{max}$, and the sign of the group velocity $c_g =
\partial\omega/\partial k$.

Experimentally, only naturally arising waves have been studied.  Darbois~Texier
\textit{et al.}~\cite{DarboisTexier2020} detected waves by laser profilometry but
did not control the perturbation wavenumber and did not measure $\sigma(k)$.
No experimental dispersion relation for oobleck waves exists.  This is a
remarkable gap given the maturity of the theory.

\begin{mdframed}[style=questionbox]
\textbf{Key questions.}
(1)~Does $\sigma(k)$ from experiment agree with the WC-based linear stability
theory \cite{DarboisTexier2023}, with no free parameters?
(2)~Is the neutral stability wavenumber $k_\text{neutral}(\phi,\theta)$
consistent with the predicted band of unstable modes?
(3)~Does the measured phase speed $c(k)$ match theory, including the sign
of $c_g$?
\end{mdframed}

\subsection{Geometry, perturbation scheme, and scaling}

A steady base flow at controlled thickness $h_0$ and mean velocity $\bar{U}$
is established in a recirculating inclined-plane channel (length $L = 80$--$150$~cm,
width $W = 10$--$15$~cm, slope $\theta$ adjustable from $0^\circ$ to $30^\circ$;
pump with PID-controlled flow rate).  A thin Teflon paddle ($\sim$2~mm height,
spanning $W$) at the inlet is driven sinusoidally at frequency $f_0$ by a
voice-coil actuator.  This imposes a surface-height perturbation of dominant
wavenumber $k_0 = 2\pi f_0/\bar{U}$ at $x = 0$.  Sweeping $f_0$ at fixed
base-flow conditions maps the full $\sigma(k)$ curve.

\subsection{Free-surface height measurement}

Darbois~Texier \textit{et al.}~\cite{DarboisTexier2020} measured $h(x,t)$
using a Keyence laser line profilometer, which resolves the opaque
suspension-air interface by triangulation.  A simpler and equally
effective approach for our setup is backlit side-view imaging: a diffuse
LED panel behind the channel provides high-contrast silhouettes of the
opaque suspension, and the high-speed camera positioned at the side records
$h(x,t)$ directly at each frame.  This method requires no additional
instrumentation beyond the existing camera.  The free-surface perturbation
$\eta(x,t) = h(x,t) - h_0$ is extracted by edge detection.

The amplitude $|\hat{\eta}(x,f_0)|$ at the forcing frequency is extracted at
multiple streamwise positions $x_1,\ldots,x_n$ by lock-in analysis, yielding
the spatial growth rate
\begin{equation}
  \sigma(k_0) = \frac{1}{\Delta x}\ln\frac{|\hat{\eta}(x_2)|}{|\hat{\eta}(x_1)|}.
  \label{eq:sigma}
\end{equation}
The phase velocity $c(k_0)$ is obtained by cross-correlating $\eta(x_1,t)$
and $\eta(x_2,t)$ to find the propagation delay.

\subsection{Practical constraints}

Cornstarch suspensions are optically opaque at $\phi \gtrsim 0.40$; a laser
sheet cannot penetrate the bulk.  The primary measurement is therefore the
free-surface profile $h(x,t)$ from side-view imaging, not interior velocity
fields.  Cross-sectional PIV inside the flow would require either (a)~an
ultrasound velocimetry system or (b)~a fully transparent index-matched
suspension (e.g., PMMA beads in water--DMSO or silica beads in a
refractive-index-matched glycerol--water mixture); both are reserved for
future work.  The surface-height measurement alone is sufficient to map
$\sigma(k)$ and $c(k)$, which are the primary deliverables of this experiment.

Sedimentation: cornstarch suspensions at $\phi < 0.43$ can sediment on the
timescale of an experiment.  The Bougouin group discards flows with $H < 7.5$~mm
for this reason~\cite{Bougouin2026PRL}; we apply the same criterion.
Recirculating the suspension through a pump continuously re-mixes it.

\begin{mdframed}[style=setupbox]
\textbf{Proposed setup.}
Recirculating channel (acrylic, $150 \times 12 \times 5$~cm, variable slope
mount); gear pump with PID flow control; inlet paddle voice-coil (0.1--50~Hz).
Backlit LED panel + side-view high-speed camera at 2000--5000~fps for $h(x,t)$
over a 60~cm window.  Cornstarch-water at $\phi = 0.40$, $0.43$, $0.45$, $0.46$;
$\theta = 5^\circ$--$25^\circ$; $h_0 = 5$--$15$~mm.  Lock-in analysis of
$\eta(x,t)$ at each $(f_0, \phi, \theta)$.
\end{mdframed}

\subsection{Observables and expected results}

Plot $\sigma$ vs.\ $k_0$ for each $(\phi, \theta)$: expect a band of unstable
wavenumbers $k_\text{min} < k_0 < k_\text{max}$ with a peak at $k_\text{max}$.
Compare directly with the Darbois~Texier (2023) predictions using
$\tauc(\phi)$ and $\etaz(\phi)$ from independent rheometry.  Plot $c(k_0)$
vs.\ $k_0$: theory predicts $c > \bar{U}$ in the oobleck-wave regime
(waves travel faster than mean flow).

\begin{mdframed}[style=noveltybox]
\textbf{Novelty and publishability.}
This is the first experimental measurement of the oobleck-wave dispersion
relation.  The controlled-perturbation protocol distinguishes stable from
unstable wavenumbers and measures growth rates without free parameters,
providing a stringent quantitative test of the WC-based stability theory that
the existing literature cannot supply.  A clean $\sigma(k)$ curve across two
or three volume fractions constitutes a full \textit{Journal of Fluid Mechanics}
or \textit{Physical Review Fluids} paper.  Extension to interior velocity
measurement (via index-matched transparent system or ultrasound velocimetry)
is a natural second phase.
\end{mdframed}

%% ---------------------------------------------------------------
\section{Experiment~4: Stairwave Wavelength Selection Under Controlled Tilt Rate}
\label{sec:stairwave}

\subsection{Physical motivation and gap}

Balmforth~\cite{Balmforth2025} showed theoretically that the nonlinear
saturation of oobleck waves produces steadily propagating, finite-amplitude
stairwaves: steep fronts separated by nearly flat regions at the base-flow
thickness.  Bougouin \textit{et al.}~\cite{Bougouin2026PRL} subsequently
observed stairwaves when a horizontal layer of DST suspension was abruptly
tilted.  The observed wavelength agrees with a geometric mass-conservation
argument:
\begin{equation}
  \lambda \approx \frac{2\,h_\text{ini}}{\sin\theta},
  \label{eq:lambda_geom}
\end{equation}
which contains no rheological parameters and predicts $\lambda$ independent
of $\phi$.

Two questions remain open.  First, in the Bougouin experiments the tilt was
rapid but not controlled; whether $\lambda$ depends on the tilt rate $\Omega$
is unknown.  Balmforth's theory connects stairwaves to the nonlinear saturation
of the fastest-growing oobleck wave mode, which has a wavenumber $k_\text{max}$
that depends on $\phi$ and $\theta$; if the tilt is slow enough that the
suspension evolves quasi-statically through the Oobleck-unstable regime, it is
plausible that the selected $\lambda$ differs from Eq.~\eqref{eq:lambda_geom}.
Whether and how this occurs is an empirical question that the theory does not
yet prescribe analytically.  Second, how quickly do stairwaves form from an
initially flat layer, and does this timescale scale with the oobleck-wave
e-folding time?

\begin{mdframed}[style=questionbox]
\textbf{Key questions.}
(1)~Does $\lambda$ depend on $\Omega$, or is Eq.~\eqref{eq:lambda_geom} exact
regardless of tilt rate?
(2)~Does $\lambda$ depend on $\phi$ at fixed $(h_\text{ini}, \theta)$?
(If yes, rheology matters; Eq.~\eqref{eq:lambda_geom} is incomplete.)
(3)~Is there a minimum $\Omega$ below which stairwaves do not form?
(4)~What is the stairwave formation time $t_\text{stair}(\phi, \theta, \Omega)$?
\end{mdframed}

\subsection{Geometry, protocol, and scaling}

A flat acrylic tray ($40 \times 15$~cm, depth 2~cm) is mounted on a motorized
hinge (stepper motor, angular resolution $0.05^\circ$, range 0.1--$90^\circ$/s)
and filled with a suspension layer of thickness $h_\text{ini}$ at rest.  The
motor tilts the tray to a target angle $\theta_f$ at rate $\Omega$.

A natural timescale for comparing the tilt to the wave dynamics is
$\Omega^* \sim \theta_f \cdot \sigma_\text{max}(\phi, \theta_f)$, where
$\sigma_\text{max}$ is the peak temporal growth rate of the oobleck instability
(dimension s$^{-1}$).  Roughly, if $\Omega \gg \Omega^*$ the tray reaches its
final angle before Oobleck waves can grow; if $\Omega \ll \Omega^*$ the waves
begin growing during the tilt transient.  From the linear stability theory of
Darbois-Texier \textit{et al.}~\cite{DarboisTexier2023}, $\sigma_\text{max}$ at
$\phi = 0.45$ and $\theta \sim 15^\circ$ is of order $10^{-2}$~s$^{-1}$, giving
$\Omega^* \sim 10^{-2} \times 15^\circ \sim 0.15^\circ$--$1^\circ$/s (the
range reflects uncertainty in the exact $\sigma_\text{max}$ without running the
full linear stability code).  The motor range (0.1--$90^\circ$/s) straddles
this estimate by over two decades in either direction, which is the essential
experimental requirement.

\begin{mdframed}[style=setupbox]
\textbf{Proposed setup.}
Motorized tray as described; $\phi \in \{0.43, 0.45, 0.46\}$;
$h_\text{ini} \in \{3, 5, 7, 10\}$~mm; $\theta_f \in \{10^\circ, 15^\circ,
20^\circ, 25^\circ\}$; $\Omega \in \{0.1, 0.5, 2, 10, 50\}^\circ$/s.
Side-view high-speed camera (2000~fps) for $h(x,t)$ via backlit LED; spatiotemporal
diagrams $h(x,t)$ reveal stairwave nucleation time and speed.
Top-view camera for plan-view stairwave geometry near sidewalls.
No laser or PIV required.
\end{mdframed}

\subsection{Observables and expected results}

Primary observable: spatiotemporal diagram $h(x,t)$ from side-view footage.
Stairwaves appear as diagonal stripes with slope $dx/dt = \ULD$; their spatial
spacing gives $\lambda$.  Plot $\lambda(\Omega)$ at fixed $(\phi, \theta_f,
h_\text{ini})$.

If Eq.~\eqref{eq:lambda_geom} is exact, $\lambda$ is independent of $\Omega$
and $\phi$---a clean null result that itself confirms the geometric picture and
the universality of the liquid-dam speed.  If $\lambda(\Omega)$ shows a trend,
the crossover rate $\Omega^*$ provides a direct empirical measurement of the
oobleck-wave growth rate without a separate perturbation experiment.  Measuring
$\lambda(\phi)$ at fixed $(\Omega, \theta_f, h_\text{ini})$ tests whether the
rheological parameters enter the wavelength selection at all.

\begin{mdframed}[style=noveltybox]
\textbf{Novelty and publishability.}
Stairwaves were first reported by Bougouin \textit{et al.}~\cite{Bougouin2026PRL};
this would be the first controlled study of their formation.  The experiment
requires only a motorized tray and existing camera, is executable within 2--3
weeks, and the cleanest outcome---$\lambda$ independent of $\Omega$ and $\phi$---is
itself publishable as a concise confirmation of the Balmforth--Bougouin picture.
A $\phi$-dependent or $\Omega$-dependent $\lambda$ is a more striking result and
would motivate a direct comparison with Balmforth's nonlinear theory.
Target: \textit{Physical Review Letters} or \textit{Physical Review Fluids}.
\end{mdframed}

%% ---------------------------------------------------------------
\section{Experiment~5: Liquid Dam Arrest and Reformation at a Step Discontinuity}
\label{sec:step}

\subsection{Physical motivation and gap}

The gravito-rheological model of Bougouin \textit{et al.}~\cite{Bougouin2026PRL}
predicts that $\ULD$ is an intrinsic material property: it depends only on
$\tauc(\phi)$, $\etaz(\phi)$, $\rho$, and $g$, with no sensitivity to geometry,
released volume, or surface topography.  This has been tested only on smooth,
flat surfaces.  A step discontinuity provides a decisive test: if $\ULD$ is
truly intrinsic, the dam should arrest at the step, undergo a transient, and
then re-establish at exactly $\ULD$ below the step.  Beyond this falsification
test, an upward step (barrier) of carefully chosen height can arrest the
frictionless overflow layer, providing the first indirect measurement of the
overflow-layer thickness $\delta$---the one quantity in the Bougouin model never
directly observed.

\begin{mdframed}[style=questionbox]
\textbf{Key questions.}
(1)~Does the liquid dam reform at $\ULD(\phi)$ after a downward step, confirming
the intrinsic character of the gravito-rheological speed?
(2)~What is the transient length $\Delta x_\text{trans}$ over which the dam
re-establishes, and how does it scale with $\ell$ and $\Delta z/H$?
(3)~What is the critical barrier height $b^*$ above which the jammed front
is permanently arrested, and does it scale as predicted from the overflow
thickness $\delta$?
(4)~Is there an upstream-propagating disturbance when the front is arrested?
\end{mdframed}

\subsection{Geometry and scaling}

\paragraph{Downward step.}
The jammed front arrives at the step with speed $\ULD$ and height $H_1$.
Below the step, the hydrostatic driving increases by $\rho g\Delta z$.  Since
$\ULD$ is set by local rheology alone (independent of $\Delta z$), the dam must
reform at the same speed below the step.  The reformation involves a transient
in which the overflow layer re-equilibrates to its original thickness $\delta$.
We do not make a detailed analytical prediction for $\Delta x_\text{trans}$---the
model of Bougouin \textit{et al.}\ is not derived for this configuration---but
parametrically, the relevant scales are $\ell$ (the jammed-front length) and
$\Delta z/\delta$ (the step height normalized by the overflow thickness).

\paragraph{Upward barrier (measurement of $\delta$).}
An upward barrier of height $b$ blocks the jammed front but allows the thin
frictionless overflow to pass over if $b < \delta$.  The critical barrier
height that just arrests the overflow is $b^* \approx \delta$.  From the model,
the overflow thickness is set by the DST onset condition on the overflow shear
stress: the shear stress at the base of the overflow equals $\tauc$, giving
\begin{equation}
  \frac{\rho g\,\delta^2}{2\ell} \sim \tauc(\phi),
  \qquad \Rightarrow \qquad
  \delta \sim \left(\frac{2\ell\,\tauc}{\rho g}\right)^{1/2},
  \label{eq:delta}
\end{equation}
and hence $b^* \sim \delta$.  With $\ell \approx 4H$ (from the Bougouin
surface velocimetry), $H = 3$~cm, $\tauc = 3$~Pa, and $\rho g = 11{,}760$~N/m$^3$,
\begin{equation}
  \delta \approx \left(\frac{2 \times 0.12\,\text{m} \times 3\,\text{Pa}}
  {11{,}760\,\text{N/m}^3}\right)^{1/2} \approx 8\,\text{mm}.
  \label{eq:delta_num}
\end{equation}
This estimate uses $\ell = 4H$ from the $\ell/H \approx 3$--$5$ range reported
by Bougouin \textit{et al.}  Steps in the range 3--15~mm comfortably bracket
$b^*$.  Equation~\eqref{eq:delta} predicts $b^* \propto (H\tauc)^{1/2}$, which
can be tested by varying $H$ and $\phi$ independently.

We note that Eq.~\eqref{eq:delta} is derived from the Bougouin model and
involves $\ell$, which is itself uncertain by a factor of $\sim 2$.  The
numerical estimate in Eq.~\eqref{eq:delta_num} should be treated as an
order-of-magnitude guide for step sizing, not a precise prediction.

\begin{mdframed}[style=setupbox]
\textbf{Proposed setup.}
Primary channel ($100 \times 12 \times 6$~cm acrylic, horizontal).  Modular
acrylic floor inserts, shimmed to precise heights, provide steps of
$\Delta z \in \{-10, -5, -3, +1, +3, +5, +8, +12\}$~mm at $x_s = 30$~cm
from the release point.  Dam-break release: acrylic wedge ($A = 50$--$200$~cm$^2$)
initially at $45^\circ$, tilted to horizontal (identical to the Bougouin
protocol).  Side-view high-speed camera (500--2000~fps) images the full
channel.  Spatiotemporal diagram $X(t)$ extracted by front-edge detection.
Surface PTV (seeded surface, horizontal laser sheet) for overflow velocity
profile before and after step.
\end{mdframed}

\subsection{Observables and expected results}

\emph{Downward step:} Spatiotemporal plot $X(t)$ shows a possible kink when
the front crosses $x_s$, followed by recovery to slope $\ULD$.  The recovery
length $\Delta x_\text{trans}$ vs.\ $\Delta z$ and $H$ characterizes the
jammed-front relaxation.

\emph{Upward barrier:} Front speed $U_\text{post}(b)$ measured for each $b$.
For $b \ll b^*$: dam passes with no speed change.  For $b \sim b^*$: dam
slows or arrests.  For $b \gg b^*$: front permanently arrested.  The threshold
$b^*(\phi, H)$ tests Eq.~\eqref{eq:delta}.  Surface PTV at the step: for
$b < b^*$, the overflow-layer velocity profile above the step directly reveals
the overflow dynamics; comparing $u_r(r)$ just upstream and downstream of the
barrier provides the most direct evidence of the overflow layer in either configuration.

\begin{mdframed}[style=noveltybox]
\textbf{Novelty and publishability.}
The step provides a decisive, falsifiable test of the intrinsic character of
$\ULD$---something no existing liquid-dam experiment has done.  The barrier
threshold $b^*$ gives the first indirect measurement of $\delta$, closing the
last major gap in the Bougouin model.  Both sub-experiments are low-cost
(acrylic shims, existing channel, existing camera) and executable within a few
weeks.  The combination of model falsification and internal-structure measurement
in a single experiment makes this a strong \textit{Physical Review Letters}
candidate.
\end{mdframed}

%% ---------------------------------------------------------------
\section{Instrumentation Notes}
\label{sec:measurement}

\subsection{High-speed photography}

All five experiments use a single side-view or top-view high-speed camera as
the primary diagnostic.  Frame rates of 500--2000~fps are sufficient for all
proposed flow velocities ($\ULD \sim 5$--$50$~mm/s) and length scales
($H \sim 5$--$30$~mm).  Backlit LED panel illumination provides excellent
contrast at the jammed front (the opacity change between flowing and jammed
cornstarch is visible to the naked eye and is stark in high-speed images).

\subsection{Laser PIV / PTV}

Experiments 1, 2, and 5 use surface PTV: a CW 532~nm DPSS laser
($\geq 300$~mW) expanded to a $\sim$1~mm horizontal sheet illuminates hollow
glass microspheres ($50$--$100$~\textmu m) floating at the free surface.  The
top-view high-speed camera (with a 550~nm long-pass filter to reject scattered
laser light) records particle tracks.  This requires no additional hardware
beyond a CW laser, cylindrical lens, and band-pass filter---all items available
from standard optics suppliers.

Cross-sectional (volumetric) PIV inside the bulk of an opaque cornstarch
suspension is not possible with visible-light laser sheets and is not proposed
here.  Any future requirement for interior velocity fields should be addressed
by (a)~acoustic velocimetry (ultrasound Doppler), (b)~an index-matched
transparent particle system, or (c)~high-resolution x-ray imaging, as
suggested by Bougouin \textit{et al.}~\cite{Bougouin2026PRL}.

\subsection{Rheological characterization}

Independent Couette rheometry ($\dot{\gamma}$-sweep at controlled rate)
provides $\tauc(\phi)$ and $\etaz(\phi)$ via WC fits.  The Han--van~Ha--Jaeger
cornstarch correction $\phi_\text{eff} = 1.136\,\phi_\text{nom}$
\cite{HanVanHaJaeger2017} is applied throughout.  Moisture content is
controlled by preparing suspensions in a humidity-controlled environment
($65 \pm 5\%$~RH) and measuring suspension density before each run.  Validating
key results in a second particle system (polystyrene beads in water, or
cornstarch in water--glycerol) confirms generality, following the Bougouin
protocol~\cite{Bougouin2026PRL}.

\subsection{Prioritization}

Table~\ref{tab:priority} summarizes the estimated setup complexity and time to
first publishable figure for each experiment.

\begin{table}[h]
\caption{\label{tab:priority}
Prioritization of proposed experiments.
Setup complexity: L = low (stock parts + machine shop),
M = medium (1--2 weeks fabrication + alignment).
Time to result: elapsed time from start to first publication-quality figure.}
\begin{ruledtabular}
\begin{tabular}{lllll}
Exp. & Topic & Complexity & Time & Primary target \\
\hline
4 & Stairwave tilt-rate     & L   & 2--3 weeks  & \textit{PRL} / \textit{PRF} \\
5 & Step discontinuity      & L   & 3--5 weeks  & \textit{PRL}               \\
1 & 3D radial dam break     & L   & 3--5 weeks  & \textit{PRL} / \textit{PRF} \\
2 & Elastic membrane        & M   & 4--8 weeks  & \textit{PRL} / \textit{SM}  \\
3 & Wave dispersion mapping & M   & 8--16 weeks & \textit{JFM} / \textit{PRF} \\
\end{tabular}
\end{ruledtabular}
\end{table}

Experiments 4, 5, and 1 can be started immediately with stock acrylic,
a motorized hinge, and the existing high-speed camera.  They share no hardware
and can run in parallel.  Experiments 2 and 3 require additional fabrication
(membrane frame and actuator; recirculating channel) but are the highest-impact
targets and should be designed concurrently with the first three.

%% ---------------------------------------------------------------
\section{Discussion}
\label{sec:discussion}

The five experiments collectively address three open questions in the
liquid-dam and oobleck-wave literature as of early 2026.

\begin{enumerate}[leftmargin=1.5em, label=(\roman*)]

\item \textit{Three-dimensionality} (Experiment~1):
All existing theory and experiment is quasi-2D.  The 3D case modifies the
height dynamics from $H \sim t^{-1/5}$ (Newtonian 2D) to the clean $H \sim t^{-2}$
during the liquid-dam phase, with a subsequent transition to the Huppert
axisymmetric regime.  The azimuthal stability of the jammed front is entirely
open.

\item \textit{Internal flow structure and model validation}
(Experiments 3, 5):
The overflow thickness $\delta$ and jammed-front length $\ell$ enter the
Bougouin model but have not been directly measured.  Experiment~5 measures
$\delta$ indirectly via the barrier threshold; Experiment~3 maps $\sigma(k)$
with no free parameters, providing the first quantitative test of the
Darbois~Texier stability theory.

\item \textit{Controlled dynamics} (Experiments 4, 2):
Stairwave formation has only been observed under abrupt tilt; controlling $\Omega$
converts an observation into a measurement.  Experiment~2 similarly converts the
impact scenario (impulsive, inertia-dominated) into a controlled quasi-static
forcing, isolating the WC nucleation mechanism.

\end{enumerate}

Any publishable result from these experiments should include validation in a
second particle system (polystyrene--water or cornstarch--glycerol) to confirm
that the observed behavior is generic to DST suspensions rather than specific
to cornstarch.

%% ---------------------------------------------------------------
\begin{acknowledgments}
This work is supported by [grant details].  The author thanks the Driscoll
Laboratory and the Cheng Laboratory (Northwestern University) for access to
equipment and discussion.
\end{acknowledgments}

%% ---------------------------------------------------------------
\begin{thebibliography}{99}

\bibitem{DarboisTexier2020}
B.~Darbois~Texier, H.~Lhuissier, Y.~Forterre, and B.~Metzger,
\textit{Surface-wave instability without inertia in shear-thickening suspensions},
Commun.\ Phys.\ \textbf{3}, 232 (2020).

\bibitem{DarboisTexier2023}
B.~Darbois-Texier, H.~Lhuissier, B.~Metzger, and Y.~Forterre,
\textit{Shear-thickening suspensions down inclines: from Kapitza to Oobleck waves},
J.\ Fluid Mech.\ \textbf{959}, A27 (2023).

\bibitem{Bougouin2026PRL}
A.~Bougouin, H.~Lhuissier, Y.~Forterre, and B.~Metzger,
\textit{Liquid Dam: A Constant-Speed Regime in Gravity-Driven
Shear-Thickening Suspensions},
Phys.\ Rev.\ Lett.\ \textbf{136}, 158202 (2026).

\bibitem{Bougouin2024PNAS}
A.~Bougouin, B.~Metzger, Y.~Forterre, P.~Boustingorry, and H.~Lhuissier,
\textit{A frictional soliton controls the resistance law of
shear-thickening suspensions in pipes},
Proc.\ Natl.\ Acad.\ Sci.\ U.S.A.\ \textbf{121}, e2321581121 (2024).

\bibitem{WyartCates2014}
M.~Wyart and M.~E.~Cates,
\textit{Discontinuous shear thickening without inertia in dense
non-Brownian suspensions},
Phys.\ Rev.\ Lett.\ \textbf{112}, 098302 (2014).

\bibitem{Balmforth2025}
N.~J.~Balmforth,
\textit{Nonlinear Oobleck waves},
J.\ Fluid Mech.\ \textbf{1002}, A34 (2025).

\bibitem{Rathee2022PNAS}
V.~Rathee, J.~Miller, D.~Blair, and J.~Urbach,
\textit{Structure of propagating high-stress fronts in a
shear-thickening suspension},
Proc.\ Natl.\ Acad.\ Sci.\ U.S.A.\ \textbf{119}, e2203795119 (2022).

\bibitem{Depoilly2024}
F.~Depoilly, S.~Millet, H.~Ben~Hadid, S.~Dagois-Bohy, and F.~Rousset,
\textit{Unifying the roll waves},
PLoS One \textbf{19}, e0310805 (2024).

\bibitem{Rocha2023JFM}
F.~Rocha, Y.~Forterre, B.~Metzger, and H.~Lhuissier,
\textit{Drag of a shear-thickening suspension on a rotating cylinder},
J.\ Fluid Mech.\ \textbf{970}, A35 (2023).

\bibitem{WaitukaitisJaeger2012}
S.~R.~Waitukaitis and H.~M.~Jaeger,
\textit{Impact-activated solidification of dense suspensions via dynamic
jamming fronts},
Nature \textbf{487}, 205 (2012).

\bibitem{HanVanHaJaeger2017}
E.~Han, N.~Van~Ha, and H.~M.~Jaeger,
\textit{Measuring the porosity and compressibility of liquid-suspended
solids using ultrasound},
Soft Matter \textbf{13}, 3506 (2017).

\bibitem{Mobaseri2025PNAS}
S.~Mobaseri \textit{et al.},
Proc.\ Natl.\ Acad.\ Sci.\ U.S.A.\ (2025).

\bibitem{Mobaseri2026arXiv}
S.~Mobaseri \textit{et al.},
arXiv:2601.12642 (2026).

\bibitem{Huppert1982}
H.~E.~Huppert,
\textit{The propagation of two-dimensional and axisymmetric viscous
gravity currents over a rigid horizontal surface},
J.\ Fluid Mech.\ \textbf{121}, 43 (1982).

\end{thebibliography}

\end{document}
